
%
% $Id: slides.tex 4228 2006-06-21 21:55:12Z jjamor $
%
%
% Compilar a .pdf con LaTeX (pdflatex)
% Es necesario instalar Beamer (paquete latex-beamer en Debian)
%

%
% Gráficos:
% Los gráficos pueden suministrarse en PNG, JPG, TIF, PDF, MPS
% Los EPS deben convertirse a PDF (usar epstopdf)
%

\documentclass{beamer}
\usetheme{Warsaw}
\usebackgroundtemplate{\includegraphics[width=\paperwidth]{format/libresoft-bg.png}}
\usepackage[english]{babel}
\usepackage[utf8]{inputenc}
\usepackage{graphics}
\usepackage{amssymb} % Simbolos matematicos
\usepackage{url}

%\definecolor{libresoftgreen}{RGB}{162,190,43}
%\definecolor{libresoftblue}{RGB}{0,98,143}

%\setbeamercolor{titlelike}{bg=libresoftgreen}

%% Metadatos del PDF.
\hypersetup{
  pdftitle={Libresoft},
  pdfauthor={Felipe Ortega},
  pdfcreator={GSyC/Libresoft},
  pdfproducer=PDFLaTeX,
  pdfsubject={Automated FLOSS projects metrics},
}
%%


\AtBeginSection[]
{
  \begin{frame}<presentation>
    \frametitle{Index}
    \tableofcontents[current]
  \end{frame}
}


\begin{document}

\title{Automated FLOSS projects metrics}
\subtitle{Master on Libre Software}
\institute{\\jfelipe@libresoft.es\\
GSyC/Libresoft}
\author{Felipe Ortega}
\date{\today}

\frame{
\maketitle
\begin{center}
\includegraphics[width=6cm]{format/gsyc-urjc}
\end{center}
}


% Si el titulo o el autor se quieren acortar para los pies de p�gina
% se pueden redefinir aqu�:
%\title{Titulo corto}
%\author{Autores abreviado}


%% LICENCIA DE REDISTRIBUCION DE LAS TRANSPAS
\frame{
~
\vspace{4cm}

\begin{flushright}
{\tiny
(cc) 2010 Felipe Ortega. \\
Some rights reserved. This document is distributed under the Creative \\
            Commons Attribution-ShareAlike 3.0 licence, available in \\
            http://creativecommons.org/licenses/by-sa/3.0/

%  Este documento (o uno muy similar) está disponible en \\
%  \url{http://gsyc.escet.urjc.es/~jjamor/}
}
\end{flushright}
}
%%

%%%%%%
%Transpas separadas por \begin{frame}
%%%%%%%%%%%%%%%%%%%%%%%%\end{frame}

\section{Introduction}

\begin{frame}
\frametitle{Why we need metrics about FLOSS projects?}
\begin{itemize}
\item Describe and account FLOSS project organization
\item Characterize evolution patterns in FLOSS project.
\item Finding out common structures in FLOSS projects
\end{itemize}
\end{frame}

\begin{frame}
 \frametitle{Overview of the analysis}
 \begin{itemize}
 \item Identify target project(s)
 \item Retrieve information from SCM platform (CVS, SVN, GIT...)
 \item Compute metrics and graphs locally
 \item Arrange results and compose report
 \end{itemize}
\end{frame}

\begin{frame}
 \frametitle{What we need?}
 \begin{itemize}
 \item Automated tools
 \item Libre software
 \item Support for data management and statistics
 \item MySQL/SQLite + CVSAnalY + GNU R
 \end{itemize}
\end{frame}

\section{CVSAnaly: retrieving data from SCMs}

\begin{frame}
 \frametitle{CVSAnalY: the tool}
 \begin{itemize}
 \item The CVSAnalY tool extracts information out of source 
code repository logs and stores it into a database.
 \item Support: CVS, SVN, GIT.
 \item Future support: Mercurial, Bazaar.
 \item Outcome: database of extracted information
 \end{itemize}
\end{frame}

\begin{frame}
\frametitle{CVSAnalY database schema}
 \begin{itemize}
 \item Data for every commit (\textit{scmlog})
 \item Data for every file recorded (\textit{files})
 \item Data for every contributor (\textit{people})
 \item Basic metrics: sloc, loc, ncomments, nfunctions...
 \item \texttt{\footnotesize http://melquiades.flossmetrics.org/wiki/doku.php?id=scm}
 \end{itemize}
\end{frame}
% 

\begin{frame}
\frametitle{Installing CVSAnalY}
 \begin{itemize}
 \item Dependencies:
 \begin{itemize}
  \item mysql-server
\item python (sqlite3, python-pysqlite2)
\item python-mysqldb
\item sloccount
\item automake
\item autoconf
\item git-core
\item pkg-config
 \end{itemize}
 \end{itemize}
\end{frame}

\begin{frame}
\frametitle{Installing CVSAnalY}
 \begin{itemize}
  \item git clone git://git.libresoft.es/git/cvsanaly
  \item git clone git://git.libresoft.es/git/repositoryhandler
  \item Install Repository Handler \texttt{\$ ./autogen.sh  \&\& make \&\& sudo make install} 
(or copy \textit{repositoryhandler} dir in \textit{/cvsanaly})
  \item Install CVSAnalY2 \texttt{\$ sudo python ./setup.py install}
  \item Check everything correctly works \texttt{\$ cvsanaly2 --help}
\item Retrieving info from remote URI and using SQLite as dbserver.
 \item \small \texttt{\$ cvsanaly2 --db-driver sqlite -d ~/db/nautilus.db \\}
\texttt{http://svn.gnome.org/svn/nautilus}
 \end{itemize}
\end{frame}

\section{GNU R: computing statistics}

\begin{frame}
\frametitle{Why using GNU R}
 \begin{itemize}
 \item Most widespread tool for statistical analysis
 \item More than 2,000 libraries already availabe.
 \item Perfect integration with DBs (MySQL, SQLite).
 \item Easy automation (scripts).
 \end{itemize}
\end{frame}

\begin{frame}
\frametitle{Installing GNU R}
 \begin{itemize}
 \item \texttt{\# apt-get install r-base r-recommended r-cran-rmysql}
 \item You can also install additional packages (from R CRAN).
 \item Download package, then: 
\texttt{\$ sudo R CMD INSTALL package.gz}
 \end{itemize}
\end{frame}

\begin{frame}
\frametitle{GNU R: receipts}
 \begin{itemize}
 \item Scripts computing statistics (numbers, graphs).
 \item Total figures, averages...
 \item Graphs for aggregated statistics
 \item Example of metric evolution (monthly number of commits)
 \item Advanced techniques made easy (season-trend decomposition).
 \item Usage: \texttt{\$RScript <receipt-name.R> <dbname.db>}
 \end{itemize}
\end{frame}

\end{document}